\subsection{Nützliche Formeln}
$$ax^2 + bx + c = 0 \Rightarrow x = \frac{-b \pm \sqrt{b^2 - 4ac}}{2a}$$
$$b^n- a^n = (b-a)(b^{n-1}+b^{n-2}\cdot a + \dots + a^{n-2}\cdot b + a^{n-1})$$
$$\binom{n}{k} = \frac{n!}{(n-k)!k!}; \quad \sqrt{a_k\cdot a_{k+1}} \leq \frac{a_k + a_{k+1}}{2}$$

\subsubsection{Bernoulli Ungleichung}
$$(1+x)^{n} \geq 1 + n*x \quad \forall n \in \mathbb{N}, x > -1$$

\subsubsection{Young'sche Ungleichung}
$$\forall \epsilon > 0, \forall x, y \in \mathbb{R}: 2 |xy| \leq \epsilon x^2 + \frac{1}{\epsilon} y^2$$

\subsubsection{Binomialsatz}
$$(x + y)^n = \sum_{k = 0}^n \binom{n}{k} x^{n-k}y^k$$

\subsection{Trigonometrischen Funktionen}

\bgroup
\def\arraystretch{1.8}
\begin{center}
\begin{tabular}{c|c|c|c|c|c|c|c|c|c}
     \(\alpha\)& 0& \(30^{\circ}\)& \(45^{\circ}\)& \(60^{\circ}\)& \(90^{\circ}\)& \(120^{\circ}\)& 150\(^{\circ}\)& 180\(^{\circ}\)
     & \(270^{\circ}\)\\
     \(\)& 0& \(\frac{\pi}{6}\)& \(\frac{\pi}{4}\)& \(\frac{\pi}{3}\)& \(\frac{\pi}{2}\)&\(\frac{2\pi}{3}\)& \(\frac{5\pi}{6}\)& \(\pi\) & \(\frac{3\pi}{2}\)\\
     \hline
     \(\sin{\alpha}\)& 0& \(\frac{1}{2}\)& \(\frac{\sqrt{2}}{2}\)& \(\frac{\sqrt{3}}{2}\)& 1& \(\frac{\sqrt{3}}{2}\)& \(\frac{1}{2}\)& 0
     & -1\\
     \hline
     \(\cos{\alpha}\)& 1& \(\frac{\sqrt{3}}{2}\)& \(\frac{\sqrt{2}}{2}\)& \(\frac{1}{2}\)& 0& \(-\frac{1}{2}\)& \(-\frac{\sqrt{3}}{2}\)& -1
     & 0\\
     \hline
     \(\tan{\alpha}\)& 0& \(\frac{\sqrt{3}}{3}\)& \(1\)& \(\sqrt{3}\)& N/A & \(-\sqrt{3}\)& \(-\frac{\sqrt{3}}{3}\)& 0
     & N/A \\
\end{tabular}
\end{center}

\begin{center}
\begin{tabular}{c|c|c|c|c|c|c|c}
     & \rotatebox{90}{\(-\alpha\)}& \rotatebox{90}{90\(^{\circ}\) \(-\alpha\)}& \rotatebox{90}{90\(^{\circ}\) \(+\alpha\)}& \rotatebox{90}{180\(^{\circ}\) \(-\alpha\)}& \rotatebox{90}{180\(^{\circ}\) \(+\alpha\)}& \rotatebox{90}{k*360\(^{\circ}\) \(-\alpha\)}&\rotatebox{90}{k*360\(^{\circ}\) \(+\alpha\)}\\
     \hline
     \(\sin\)&
     \(-\sin{\alpha}\)& \(\cos{ }\)& \(\cos{ }\)& \(\sin{ }\)&
     \(-\sin{ }\)& 
     \(-\sin{ }\)&
     \(\sin{ }\)\\
     \hline
     \(\cos\)&
     \(\cos{ }\)&
     \(\sin{ }\)&
     \(-\sin{ }\)&
     \(-\cos{ }\)&
     \(-\cos{ }\)&
     \(\cos{ }\)&
     \(\cos{ }\)\\
     \hline
     \(\tan\)&
     \(-\tan{ }\)&
     \(\cot{ }\)&
     \(-\cot{ }\)&
     \(-\tan{ }\)&
     \(\tan{ }\)&
     \(-\tan{ }\)&
     \(\tan{ }\)\\
\end{tabular}
\end{center}
\egroup

%----------------------------------------------------------------------------------------------------

\subsection{Bekannte Taylorreihen}

\paragraph{Gültig für alle $x\in\mathbb{R}$}
\begin{align*}
        \mathrm{e}^x &= \sum_{n=0}^{\infty}\frac{x^n}{n!}
        = 1+x+\frac{x^2}{2}+\frac{x^3}{3!}+\frac{x^4}{4!}+\mathcal{O}(x^5) \\
        \sin{x} &= \sum_{n=0}^{\infty}(-1)^n\frac{x^{2n+1}}{(2n+1)!}
        = x-\frac{x^3}{3!}+\frac{x^5}{5!}+\mathcal{O}(x^7) \\
        \sinh(x) &= \sum_{n=0}^{\infty}\frac{x^{2n+1}}{(2n+1)!}
        = x+\frac{x^3}{3!}+\frac{x^5}{5!}+\mathcal{O}(x^7) \\
        \cos(x) &= \sum_{n=0}^{\infty}(-1)^n\frac{x^{2n}}{(2n)!}
        = 1-\frac{x^2}{2}+\frac{x^4}{4!}-\frac{x^6}{6!}+\mathcal{O}(x^8) \\
        \cosh(x) &= \sum_{n=0}^{\infty}\frac{x^{2n}}{(2n)!}
        = 1+\frac{x^2}{2}+\frac{x^4}{4!}+\frac{x^6}{6!}+\mathcal{O}(x^8)
\end{align*}

\paragraph{Gültig für $|x|<1$}
\begin{align*}
        \log(1+x) &= \sum_{n=1}^{\infty}(-1)^{n+1}\frac{x^n}{n}
        = x-\frac{x^2}{2}+\frac{x^3}{3}-\frac{x^4}{4}+\mathcal{O}(x^5) \\
        (1+x)^\alpha &= \sum_{n=0}^{\infty}\binom{\alpha}{n}x^n
        = 1+\alpha x+\frac{\alpha(\alpha-1)}{2!}x^2+ \frac{\alpha(\alpha - 1)(\alpha - 2)}{3!}x^3 + \mathcal{O}(x^4) \\
        \sqrt{1+x} &= \sum_{n=0}^{\infty}\binom{1/2}{n}x^n
        = 1 + \frac{x}{2} - \frac{x^2}{8} + \frac{x^3}{16} - \mathcal{O}(x^4)
\end{align*}

\paragraph{Gültig für $|x|<\frac{\pi}{2}$}
\begin{align*}
        \tan(x) &= \sum_{n=1}^{\infty}\frac{(-1)^{n-1}2^{2n}(2^{2n}-1)B_{2n}}{(2n)!}x^{2n-1}
        = x+\frac{x^3}{3}+\frac{2x^5}{15}+\mathcal{O}(x^7) \\
        \tanh(x) &= \sum_{n=1}^{\infty}\frac{2^{2n}(2^{2n}-1)B_{2n}}{(2n)!}x^{2n-1}
        = x-\frac{x^3}{3}+\frac{2x^5}{15}+\mathcal{O}(x^7)
\end{align*}

%----------------------------------------------------------------------------------------------------

\subsection{Typische Grenzwerte / Folgen}

\bgroup
\def\arraystretch{2}
\begin{center}
	\begin{tabular}{ | l | l | }
		 \hline
                $\lim_{x\to\infty} \frac{1}{x} = 0$ & $\lim_{x\to\infty} 1 + \frac{1}{x} = 1$ \\
		 \hline
                $\lim_{x \to \infty} e^x = \infty$ & $\lim_{x \to - \infty} e^x = 0$ \\
                 \hline
                $\lim_{x \to \infty} e^{-x} = 0$ & $\lim_{x \to - \infty} e^{-x} = \infty$ \\
                 \hline
                $\lim_{x \to \infty} \frac{e^x}{x^m} = \infty$ & $\lim_{x \to - \infty} xe^x = 0$ \\
                 \hline
                $\lim_{x \to \infty} \ln(x) = \infty$ & $\lim_{x \to 0} \ln(x) = - \infty$ \\
                 \hline
                $\lim_{x \to \infty} (1 + x)^{\frac{1}{x}} = 1$ & $\lim_{x \to 0} (1 + x)^{\frac{1}{x}} = e$ \\
                 \hline
	 	$\lim_{x\to\infty} \left(1 + \frac{1}{x}\right)^b = 1$ & $\lim_{x\to\infty} \left(1 + \frac{1}{x}\right)^b = 1$ \\
		 \hline  
	 	$\lim_{x\to\infty} x^aq^x = 0, \; \forall 0 \leq q < 1$ & $\lim_{x\to\infty} n^\frac{1}{n} = 1$ \\
		 \hline
                $\lim_{x\to \pm \infty} \left(1+\frac{1}{x}\right)^x = \mathrm{e}$ & $\lim_{x\to \infty} \left(1-\frac{1}{x}\right)^x = \frac{1}{\mathrm{e}}$ \\
		 \hline
                $\lim_{x \to \pm \infty} \left(1+\frac{k}{x}\right)^{mx} = e^{km}$ & $\lim_{x \to 0} \frac{\sin{x}}{x} = 1$ \\
                 \hline
                $\lim_{x\to 0} \frac{1}{\cos(x)} = 1$ & $\lim_{x\to 0} \frac{\cos{x}-1}{x} = 0$ \\
                 \hline
                $\lim_{x\to 0} \frac{\log{1-x}}{x} = -1$ & $\lim_{x\to 0} x\log{x} = 0$ \\
                 \hline
                $\lim_{x\to 0} \frac{1-\cos{x}}{x^2} = \frac{1}{2}$ & $\lim_{x\to 0} \frac{\mathrm{e}^x-1}{x} = 1$ \\
                 \hline
                $\lim_{x\to 0} \frac{x}{\arctan{x}} = 1$ & $\lim_{x\to\infty} \arctan{x} = \frac{\pi}{2}$ \\
                 \hline
                $\lim_{x \to \infty} \left(\frac{x}{x + k}\right)^x = e^{-k}$ & $\lim_{x \to 0} \frac{e^x - 1}{x} = 1$ \\
                 \hline
                $\lim_{x \to 0} \frac{a^x - 1}{x} = \ln(a) \; \forall a > 0$ & $\lim_{x \to 0} \frac{e^{ax} - 1}{x} = a$ \\
                 \hline
                $\lim_{x \to 0} \frac{\ln(x + 1)}{x} = 1$ & $\lim_{x \to 1} \frac{\ln(x)}{x - 1} = 1$ \\
                 \hline
                $\lim_{x \to \infty} \frac{\ln(x)}{x} = 0$ & $\lim_{x \to \infty} \frac{\log(x)}{x^a} = 0$ \\
                 \hline
                $\lim_{x \to \infty} \sqrt[x]{x} = 1$ & $\lim_{x \to \infty} \frac{2x}{2^x} = 0$ \\
                 \hline 
                $\lim_{x \to \frac{\pi^-}{2}} \tan{x} = +\infty$ & $\lim_{x \to \frac{\pi^+}{2}} \tan{x} = -\infty$ \\
                \hline
                $\lim_{x \to \infty} \frac{\sin{x}}{x} = 0$ &  $\lim_{x \to 0^+} x\ln{x} = 0$\\
                \hline  
	\end{tabular}
\end{center}
\egroup

\subsection{Typische Ableitungen und Stammfunktionen}

\renewcommand\arraystretch{1.9}

\begin{center}
        \begin{tabular}{c||c}
                $F(x)$ & $F'(x) = f(x)$ \\
                \hline \hline
                
                $ c $ & $ 0 $ \\
	        $ x^a $ & $ a \cdot x^{a - 1} $ \\
                $ \frac{1}{a+1} x^{a+1} $ & $ x^a $ \\
                $ \frac{1}{a\cdot(n+1)} (ax+b)^{n+1} $ & $ (ax+b)^n $ \\
	        $ \frac{x^{\alpha+1}}{\alpha + 1} $ & $ x^\alpha, \alpha \neq -1 $ \\
                $ \sqrt x $ & $ \frac{1}{2 \sqrt x} $ \\
                $ \sqrt[n] x $ & $ \frac{1}{n} {x}^{ \frac{1}{n} -1 } $ \\
                $ \frac{2}{3} x^{ \frac{3}{2} } $ & $ \sqrt x $ \\
                $ \frac{n}{n+1} x^{ \frac{1}{n}+1 } $ & $ \sqrt[n] x $ \\
	        $ e^x $ & $ e^x $ \\
	        $ \ln(|x|) $ & $ \frac{1}{x} $ \\
                $ \log_a |x| $  &  $ \frac{1}{x \ln a} = \log_a(e) \frac{1}{x} $ \\
	        $ \sin(x) $ & $ \cos(x) $ \\
	        $ \cos(x) $ & $ -\sin(x) $ \\
	        $ \tan(x) $ & $ \frac{1}{\cos^2(x)} = 1 + \tan^2(x) $ \\
	        $ \cot(x) $ & $ \frac{1}{-\sin^2(x)} $ \\
	        $ \arcsin(x) $ & $ \frac{1}{\sqrt{1 - x^2}} $ \\
	        $ \arccos(x) $ & $ \frac{-1}{\sqrt{1 - x^2}} $ \\
	        $ \arctan(x) $ & $ \frac{1}{1 + x^2} $ \\
	        $ \sinh(x) $ & $ \cosh(x) $ \\
	        $ \cosh(x) $ & $ \sinh(x) $ \\
                $ \tanh(x) $ & $ \frac{1}{\cosh^2(x)} = 1 - \tanh^2(x) $ \\
	        $ \text{arcsinh}(x) $ & $ \frac{1}{\sqrt{1 + x^2}} $ \\
	        $ \text{arccosh}(x) $ & $ \frac{1}{\sqrt{x^2 - 1}} $ \\
	        $ \text{arctanh}(x) $ & $ \frac{1}{1 - x^2} $ \\
                $ \frac{1}{f(x)} $ & $ \frac{-f'(x)}{(f(x))^2} $ \\
                $ a^{cx} $ & $ a^{cx} \cdot c \ln a $ \\
                $ x^x $ & $ x^x \cdot (1+\ln x) \quad \scriptstyle x > 0 $ \\
                $ (x^x)^x $ & $ (x^x)^x(x+2x\ln(x)) \quad \scriptstyle x > 0 $ \\
                $ x^{(x^x)} $ & $ x^{(x^x)}(x^{x-1}+\ln x \cdot x^x (1+\ln x)) $ \\
        \end{tabular}

        \begin{tabular}{c||c}
                $F(x)$ & $F'(x) = f(x)$ \\
                \hline \hline
                
                $ \frac{1}{a} \ln (ax+b) $  &  $ \frac{1}{ax+b}$ \\
                $ \frac{ax}{c}-\frac{ad-bc}{c^2} \ln |cx+d| $  &  $ \frac{ax+b}{cx+d} $ \\
                $ \frac{1}{2a} \ln { \big| \frac{x-a}{x+a} \big| } $  &  $ \frac{1}{x^2-a^2} $ \\
                $ \frac{x}{2} f(x) + \frac{a^2}{2} \ln ( x+f(x) ) $  &  $ \sqrt{a^2+x^2} $ \\
                $ \frac{x}{2} \sqrt{a^2-x^2} - \frac{a^2}{2} \arcsin \frac{x}{|a|} $  &  $ \sqrt{a^2-x^2} $ \\
                $ \frac{x}{2} f(x) - \frac{a^2}{2} \ln {( x+f(x) )} $  &  $ \sqrt{x^2-a^2} $ \\
                $ \ln(x+\sqrt{x^2 \pm a^2}) $  &  $ \frac{1}{\sqrt{x^2 \pm  a^2}} $ \\
                $ \arcsin( \frac{x}{|a|}) $  &  $ \frac{1}{\sqrt{a^2-x^2}} $ \\
                $ \frac{1}{a} \arctan(\frac{x}{a}) $  &  $ \frac{1}{x^2+a^2} $ \\
                $ -\frac{1}{a}\cos(ax+b) $  &  $ \sin(ax+b) $ \\
                $ \frac{1}{a}\sin(ax+b) $  &  $ \cos(ax+b) $ \\
                $ -\ln|\cos(x)| $  &  $ \tan(x) $ \\
                $ \ln |\sin(x)| $  &  $ \cot(x) $ \\
                $ \ln \left|\tan(\frac{x}{2})\right| $  &  $ \frac{1}{\sin(x)} $ \\
                $ \ln \left|\tan(\frac{x}{2}+\frac{\pi}{4})\right| $  &  $ \frac{1}{\cos(x)} $ \\
                $ \frac{1}{2} (x-\sin(x)\cos(x)) $  &  $ \sin^2(x) $ \\
                $ \frac{1}{2} (x+\sin(x)\cos(x)) $  &  $ \cos^2(x) $ \\
                $ \tan(x)-x $  &  $ \tan^2(x) $ \\
                $ -\cot(x)-x $  &  $ \cot^2(x) $ \\
                $ x \arcsin(x) + \sqrt{1-x^2} $  &  $ \arcsin(x) $ \\
                $ x \arccos(x) - \sqrt{1-x^2} $  &  $ \arccos(x) $ \\
                $ x \arctan(x) - \frac{1}{2} \ln (1+x^2) $  &  $ \arctan(x) $ \\
                $ \ln(\cosh(x)) $  &  $ \tanh(x) $ \\
                $ \ln {| f(x) |} $   &   $ \frac{f'(x)}{f(x)} $ \\
                $ x \cdot (\ln |x| - 1) $  &  $ \ln |x| $ \\
                $ \frac{1}{n+1} (\ln x)^{n+1} \quad \scriptstyle  n \neq -1 $  &  $ \frac{1}{x}(\ln x)^n $ \\
                $ \frac{1}{2n} (\ln x^n)^{2} \quad \scriptstyle  n \neq 0 $  &  $ \frac{1}{x}\ln x^n $ \\
                $ \ln |\ln x| \quad \scriptstyle x>0,x\neq 1 $  &  $ \frac{1}{x \ln x} $ \\
                $ \frac{1}{b \ln a} a^{bx} $  &  $ a^{bx} $ \\
                $ \frac{cx-1}{c^2} \cdot e^{cx} $  &  $ x\cdot e^{cx} $ \\
                $ \frac{x^{n+1}}{n+1} \left( \ln x - \frac{1}{n+1} \right) \quad \scriptstyle  n \neq -1 $  &  $ x^n \ln x $ \\
                $ \frac{e^{cx} \left( c \sin (ax+b)-a \cos(ax+b) \right)}{a^2+c^2}  $  &  $ \scriptstyle e^{cx} \sin (ax+b)  $ \\
                $ \frac{e^{cx} \left( c \cos (ax+b)+a \sin(ax+b) \right)}{a^2+c^2}  $ &  $ \scriptstyle e^{cx} \cos (ax+b)  $ \\
                $ \frac{\sin^2(x)}{2} = -\frac{\cos^2(x)}{2}  $ &  $ \sin(x)\cos(x) $ \\
        \end{tabular}
\end{center}

\subsection{Diverses}

$\int \operatorname{arsinh}(x)\,dx = x \cdot \operatorname{arsinh}(x) - \sqrt{x^2 + 1} + C$

$\int \operatorname{arcosh}(x)\,dx = x \cdot \operatorname{arcosh}(x) - \sqrt{x^2 - 1} + C$

TODO: Physikalische Formeln?

$lim_{n \rightarrow \infty}(1 + \frac{x}{n})^n = e^x$
